% ---------------------------------------------------------------------
\section{重到轻形状因子与无粲 $B$ 衰变的 QCD 因子化}\label{sec:heavy-qcdf}

两个末态介子都很轻时（$\bar B\to\pi\pi$、$\pi K$ 等），接收旁观反夸克的介子 $M_1$ 也是高能的：旁观者要么以趋于零的动量分数进入 $M_1$（软重叠），要么经虚度 $\sim m_b\Lqcd$ 的硬共线胶子获得大动量（硬旁观者散射，\cref{fig:heavy-qcdf}(b)），于是出现中间标度 $\sqrt{m_b\Lqcd}$。

%%INSERT fig_qcdf.tex

\begin{proposition}[重到轻形状因子：软重叠与硬旁观者 \tagA]\label{prop:heavy-btopi-ff}
设 $\bar B\to\pi$ 跃迁中 $\pi$ 在 $\bar B$ 静止系中的能量 $E=\frac{m_B^2-q^2}{2m_B}\sim m_b/2\gg\Lqcd$（大反冲，略去 $m_\pi$）。则在 $\Lqcd/m_b$ 的领头幂次、对 $\as$ 的所有阶，三个形状因子 $F_i$（$i=+,0,T$，定义同\cite{Beneke:2003pa}）可写成
\begin{equation}\label{eq:heavy-btopi-ff}
F_i^{B\to\pi}(E)=C_i(E,\mu)\,\xi_\pi(E,\mu)+\int_0^\infty\!\dd\omega\int_0^1\!\dd u\;T_i(E;\omega,u,\mu)\,\phi_B^+(\omega,\mu)\,\phi_\pi(u,\mu),
\end{equation}
其中 (a)~$C_i$、$T_i$ 由标度 $m_b$ 与 $\sqrt{m_b\Lqcd}$ 上的短程物理决定，可按 $\as$ 展开，$C_i$ 的树图值为运动学因子，$T_i=\order{\as}$ 是硬系数与硬共线喷注函数的卷积；(b)~“软重叠”（soft overlap）形状因子 $\xi_\pi$ 与 $i$ 无关，在 SCET 中定义为同时包含硬共线、共线与软效应的矩阵元；(c)~对 $\omega$、$u$ 的积分收敛。$\phi_B^+$ 是 $B$ 介子的领头光锥分布振幅（$\omega$ 为旁观者光锥动量，$\int_0^\infty\frac{\dd\omega}{\omega}\phi_B^+=\lambda_B^{-1}$\cite{Grozin:1996pq}，演化见\cite{Lange:2003ff}），$\phi_\pi$ 见\cref{def:excl-da}。两项都 $\sim(\Lqcd/m_b)^{3/2}$。
非微扰地位：\cref{asm:intro-bridge}（MA1、MA2）。
% errtype: E-S
\uses{def:excl-da, thm:scet-decoupling, asm:intro-bridge}
\end{proposition}

\begin{proofsketch}
(1)~两步匹配 QCD $\to$ SCET$_{\rm I}$ $\to$ SCET$_{\rm II}$，依次积出硬模式与硬共线模式；在 SCET$_{\rm I}$ 中 $B$ 的软组分是超软的，经 BPS 场重定义退耦（\cref{thm:scet-decoupling}）。(2)~只含一个硬共线夸克场的重到轻流，因 $\sla vh_v=h_v$、$\sla n_A\xi=0$，对不同 Dirac 结构只差运动学因子，给出普适的 $\xi_\pi$；多含一个硬共线胶子场的流在积出硬共线模式后化为喷注函数与 $\phi_B^+$、$\phi_\pi$ 的卷积（E-S 型陈述）。(3)~前一类矩阵元写成卷积会在端点发散；Beneke 与 Feldmann 把这些效应吸收进 $\xi_\pi$ 的定义，证明了\cref{eq:heavy-btopi-ff}对所有阶成立且卷积收敛\cite{Beneke:2003pa}；BPS 用另一种标度分离得到无端点奇异性的硬核\cite{Bauer:2002aj}（另见\cite{Hill:2002vw,Beneke:2005gs}）。
\end{proofsketch}

\paragraph{地位与缺口}
全阶论证（[A]），前提是\cref{rem:scet-assumptions}的 SCET 前提（模式完备与零箱减除，E-S 型陈述）。被因子化的只是“软重叠 + 硬旁观者”的分解：$\xi_\pi$ 不能再写成分布振幅的卷积，否则端点发散（同类现象见\cref{prop:heavy-qcdf-endpoint}与\cref{prop:excl-endpoint}，[P]）；对按非相对论处理的 $B_c\to\eta_c$，软重叠中双对数的全阶结构由一组尚无闭式解的耦合积分方程描述\cite{Bell:2024bxg}。升级途径：对 $B\to\pi$ 以积分形式证明模式完备。

\paragraph{记号与幂次计数。}
以 $b\to s$ 为例，$\mathcal H_{\rm eff}=\frac{G_F}{\sqrt2}\sum_{p=u,c}V_{pb}V_{ps}^*(C_1Q_1^p+C_2Q_2^p+\cdots)$，$Q_1^p=(\bar pb)_{V-A}(\bar sp)_{V-A}$，$Q_2^p=(\bar p_ib_j)_{V-A}(\bar s_jp_i)_{V-A}$，省略号为企鹅算符；NDR 方案中 $\mu=m_b=4.2$\,GeV 处 $C_1=1.081$、$C_2=-0.190$\cite[表~1]{BenekeEtAl2001}。领头扭度分布振幅 $\Phi_M(x)=6x\bar x[1+\sum_{n\ge1}\alpha_n^MC_n^{(3/2)}(2x-1)]$（$\alpha_n^M$ 即\cref{def:excl-da}中的 $a_n$），扭度 3 振幅取 $\Phi_p=1$，$r_\chi^M=2\mu_M/m_b$（如 $\mu_\pi=m_\pi^2/(m_u+m_d)$），$\int_0^1\frac{\dd\xi}{\xi}\Phi_B=\frac{m_B}{\lambda_B}$。$M_1$ 接收旁观者，$M_2$ 由弱顶点发射；振幅写成 $A_{M_1M_2}=\ii\frac{G_F}{\sqrt2}m_B^2F_0^{B\to M_1}(0)f_{M_2}$（朴素因子化矩阵元\cite[式~(11)]{BenekeNeubert2003}）乘以系数 $a_i$。由 $F^{B\to\pi}(0)\sim(\Lqcd/m_b)^{3/2}$、$f_B\sim\Lqcd^{3/2}m_b^{-1/2}$、$f_\pi\sim\lambda_B\sim\Lqcd$，形状因子项 $\sim m_b^2F^{B\to\pi}(0)f_\pi$ 与硬旁观者项 $\sim f_Bf_{M_1}f_{M_2}m_b/\lambda_B$ 同为 $\Lqcd^{5/2}m_b^{1/2}$。旁观者（动量 $\xi p_B$，$\xi\sim\Lqcd/m_b$）以分数 $\bar y=1-y$ 进入 $M_1$ 时，交换胶子的动量 $k=\xi p_B-\bar yp_1$ 满足（用 $2p_B\cdot p_1=m_B^2$）
\begin{equation}\label{eq:heavy-hc-gluon}
k^2=\xi^2m_B^2-\bar y\,\xi\,m_B^2\sim-m_b\Lqcd\qquad(\bar y=\order{1}),
\end{equation}
是硬共线的；BBNS 因此在 $\mu_h=\sqrt{\Lambda_h\mu}$（$\Lambda_h=0.5$\,GeV）处取硬旁观者项的 $\as$\cite{BenekeEtAl2001}。

\begin{proposition}[无粲两体衰变的领头幂次 QCD 因子化 \tagP]\label{prop:heavy-qcdf-lp}
设 $M_1$、$M_2$ 为味非单态的轻赝标介子。BBNS 因子化公式
\begin{equation}\label{eq:heavy-qcdf}
\begin{split}
\langle M_1M_2|Q_i|\bar B\rangle={}&F^{B\to M_1}(0)\int_0^1\!\dd x\,T^{\rm I}_i(x)\,f_{M_2}\Phi_{M_2}(x)\\
&+\int_0^1\!\dd\xi\,\dd x\,\dd y\;T^{\rm II}_i(\xi,x,y)\,f_B\Phi_B(\xi)\,f_{M_1}\Phi_{M_1}(y)\,f_{M_2}\Phi_{M_2}(x)
\end{split}
\end{equation}
（旁观者可进入任一介子时再加 $M_1\leftrightarrow M_2$ 的项；归一化使树图的 $T^{\rm I}_i$ 恰为朴素因子化，$T^{\rm II}_i=\order{\as}$）作为对 $\as$ 所有阶的领头幂次陈述没有公认的证明；本命题断言它在有限阶已被确立的内容。
(a)~（单圈）在 $\order{\as}$，“非因子化”的顶点修正与旁观者图中的软、共线发散在领头幂次相消，$F^{B\to M_1}$ 按\cref{prop:heavy-btopi-ff}理解；顶点修正给出
\begin{gather}
a_{1,\rm I}=C_1+\frac{C_2}{N_c}\Big[1+\frac{C_F\as}{4\pi}V_{M_2}\Big],\qquad a_{2,\rm I}=C_2+\frac{C_1}{N_c}\Big[1+\frac{C_F\as}{4\pi}V_{M_2}\Big],\label{eq:heavy-ai}\\
V_M=12\ln\frac{m_b}{\mu}-18+\int_0^1\!\dd x\,g(x)\,\Phi_M(x)\quad(\text{NDR}),\label{eq:heavy-VM}\\
\int_0^1\!\dd x\,g(x)\Phi_M(x)=-\frac12-3\ii\pi+\Big(\frac{11}{2}-3\ii\pi\Big)\alpha_1^M-\frac{21}{20}\alpha_2^M+\cdots,\label{eq:heavy-gegen}
\end{gather}
% VF-BQ-1, VF-BQ-2
其中 $g(x)=3\big(\frac{1-2x}{1-x}\ln x-\ii\pi\big)+\big[2\,\mathrm{Li}_2(x)-\ln^2x+\frac{2\ln x}{1-x}-(3+2\ii\pi)\ln x-(x\leftrightarrow1-x)\big]$（$x$ 的约定同\cite{BenekeEtAl2001}），系数 $12$ 由 $a_{1,2}$ 与 $\mu$ 无关唯一确定（对任意 $N_c$）。
(b)~（高阶）$T^{\rm I}$ 的红外有限性已显式验证到两圈（NNLO），$T^{\rm II}$ 的到其单圈修正（$\order{\as^2}$）。
非微扰地位：\cref{asm:intro-bridge}（MA1、MA2）。
% errtype: E-S
\uses{prop:heavy-btopi-ff, def:excl-da, asm:intro-bridge}
\end{proposition}

\begin{proofsketch}
(a)~是 BBNS 的单圈计算\cite[式~(46)--(48)]{BenekeEtAl2001}。系数 $12$ 可独立导出：设 $\frac{\dd\bm C}{\dd\ln\mu}=\frac{\as}{4\pi}\gamma^{\tp}\bm C$，$\gamma=\begin{psmallmatrix}-6/N_c&6\\6&-6/N_c\end{psmallmatrix}$（$C_\pm=C_1\pm C_2$ 的本征值 $\pm6(N_c\mp1)/N_c$；$\mu=m_b$ 处 $C_+=0.891<1<C_-=1.271$，与其符号相符），则
\begin{equation}\label{eq:heavy-a1-rg}
\frac{\dd}{\dd\ln\mu}\Big(C_1+\frac{C_2}{N_c}\Big)=\frac{\as}{4\pi}\Big[\Big(-\frac{6}{N_c}C_1+6C_2\Big)+\frac1{N_c}\Big(6C_1-\frac6{N_c}C_2\Big)\Big]=\frac{\as}{4\pi}\,\frac{12C_F}{N_c}\,C_2 ,
\end{equation}
恰被 $\frac{C_2}{N_c}\frac{C_F\as}{4\pi}12\ln\frac{m_b}{\mu}$ 的导数抵消（准确到 $\order{\as}$）；$a_{2,\rm I}$ 同理，常数 $-18$ 依赖于方案。\cref{eq:heavy-gegen}的系数由数值积分重现到 $10^{-20}$。数值上，取渐近分布振幅、$\mu=m_b$ 与 $\as(m_b)=0.223$，$V_\pi=-18.5-3\ii\pi$，$a_{1,\rm I}\approx1.045+0.014\,\ii$，$a_{2,\rm I}\approx0.012-0.080\,\ii$：领头幂次的强相位是 $\order{\as}$ 的微扰量（来自轻夸克中间态的在壳部分），色压低的 $a_2$ 中修正与树图值 $C_2+C_1/N_c=0.170$ 同量级，且那里硬旁观者贡献占主导\cite[第~4.2~节]{BenekeEtAl2001}（脚本 \texttt{s2\_\allowbreak qcdf.py}）。(b)~两圈顶点\cite{Bell:2007tv,Bell:2009nk,Beneke:2009ek}与企鹅振幅\cite{Bell:2015koa,Bell:2020qus}，单圈旁观者散射\cite{Beneke:2005vv,Beneke:2006mk,Kivel:2006xc,Pilipp:2007mg,Beneke:2005gs}：这些硬散射核都没有残留的软或共线发散。
\end{proofsketch}

\paragraph{地位与缺口}
部分（[P]）。一般断言（BBNS 公式在领头幂次对 $\as$ 所有阶成立）没有公认的证明，BBNS 本人指出须把 $B\to D\pi$ 的两圈证明推广到两个轻介子的末态\cite[第~8~节]{BenekeEtAl2000}。SCET 中有领头幂次的全阶推导：Chay 与 Kim 论证振幅对 $\as$ 所有阶因子化\cite{Chay:2003ju}；Bauer、Pirjol、Rothstein、Stewart（BPRS）认为其定理在 $\as(E_M\Lqcd)$ 不可微扰时也成立，并认为粲夸克企鹅（charming penguin）可能在领头阶给出长程效应\cite{Bauer:2004tj}；BBNS 认为二者的公式等价\cite{Beneke:2004bn}，BPRS 不同意其批评\cite{Bauer:2005wb}。本书不作裁断；分歧集中在硬旁观者项能否按 $\as(\sqrt{m_b\Lqcd})$ 展开与粲企鹅的幂次计数（E-S 型陈述）。因此标签取 [P] 而非 [A]；但说“BBNS 只在 $\order{\as}$ 成立”则低估了现状。升级途径：包含这两点、被普遍接受的全阶证明。

\begin{proposition}[幂次修正中的端点发散 \tagV]\label{prop:heavy-qcdf-endpoint}
在 $\order{\as}$ 按标准共线因子化计算 $\bar B\to M_1M_2$ 的扭度 3 手征增强硬旁观者项与弱湮灭项（$\Lqcd/m_b$ 幂次修正；$r_\chi^M$ 形式上 $\sim\Lqcd/m_b$，数值上 $\order{1}$）。对端点行为一般的分布振幅（扭度 2 振幅在端点线性趋于零，扭度 3 振幅在端点不为零），尤其对渐近振幅，卷积对数发散，分别含
\begin{equation}\label{eq:heavy-XAXH}
X_H=\int_0^1\frac{\dd y}{1-y}\,\Phi_p(y),\qquad X_A=\int_0^1\frac{\dd y}{y},
\end{equation}
% VF-BQ-3
故这些修正的标准共线因子化不成立。BBNS 的参数化 $X_{A,H}=(1+\varrho_{A,H}\ee^{\ii\varphi_{A,H}})\ln\frac{m_B}{\Lambda_h}$（$\varrho_{A,H}\le1$，相位任意，$\Lambda_h=0.5$\,GeV）是模型，不是因子化定理；例如对两个赝标介子、按 Beneke--Neubert 的规则 $\int_0^1\frac{\dd y}{y}\to X_A$\cite[式~(56)]{BenekeNeubert2003}，
\begin{equation}\label{eq:heavy-annih}
A_1^i\approx\pi\as\Big[18\Big(X_A-4+\frac{\pi^2}{3}\Big)+2r_\chi^2X_A^2\Big],\qquad A_3^f\approx12\pi\as\,r_\chi\big(2X_A^2-X_A\big).
\end{equation}
% VF-BQ-4, VF-BQ-5
（本命题是关于硬散射核与分布振幅之卷积的陈述。）
\uses{prop:heavy-qcdf-lp}
\end{proposition}

\begin{proof}[发散性的显式计算]
\emph{硬旁观者项。}$T^{\rm II}$ 给出（以 $\Phi_{M_2}(x)/\bar x$ 加权并对 $x$ 积分的）$\int_0^1\frac{\dd y}{\bar y}[\Phi_{M_1}(y)+r_\chi^{M_1}\frac{\bar x}{x}\Phi_p(y)]$\cite[式~(57)]{BenekeEtAl2001}。扭度 2 部分有限（渐近振幅给出 $\int_0^16y\,\dd y=3$）；扭度 3 部分为 $r_\chi^{M_1}\langle x^{-1}\rangle_{M_2}X_H$，而 $\int_0^{1-\delta}\frac{\dd y}{1-y}\Phi_p(y)=\Phi_p(1)\ln\frac1\delta+\order{1}$。发散来自旁观者作为软夸克进入 $M_1$ 的区域 $\bar y\to0$，此时由\cref{eq:heavy-hc-gluon} $k^2\to\xi^2m_B^2\sim\Lqcd^2$，交换胶子不再是硬共线的；物理上它在 $\bar y\sim\Lqcd/m_b$ 处被截断，$X_H\sim\ln(m_b/\Lqcd)$，系数可以是复数\cite{BenekeEtAl2001}。
\emph{弱湮灭项。}$A_1^i=\pi\as\int\!\!\int\dd x\,\dd y\,\{\Phi_{M_2}(x)\Phi_{M_1}(y)[\frac{1}{y(1-x\bar y)}+\frac{1}{\bar x^2y}]+r_\chi^{M_1}r_\chi^{M_2}\frac{2}{\bar xy}\}$\cite[式~(62)]{BenekeEtAl2001}。取渐近振幅，第一项有限，$36\int\!\!\int\frac{x\bar x\bar y}{1-x\bar y}=18(\frac{\pi^2}{3}-3)$；第二项 $=18\int_0^1\frac{x\,\dd x}{1-x}$ 与第三项 $=2r_\chi^2\int_0^1\frac{\dd x}{1-x}\int_0^1\frac{\dd y}{y}$ 发散（一般的扭度 2 振幅在 $x\to1$ 处同样给出 $\Phi_{M_2}/\bar x^2\propto1/\bar x$）。把 $\int_0^1\frac{\dd y}{y}$、$\int_0^1\frac{\dd x}{1-x}$ 都换成 $X_A$，用 $\frac{x}{1-x}=\frac1{1-x}-1$，三项为 $18(\frac{\pi^2}{3}-3)$、$18(X_A-1)$、$2r_\chi^2X_A^2$，相加即\cref{eq:heavy-annih}。$A_3^f$ 的两项\cite[式~(62)]{BenekeEtAl2001}在用 $\frac{x(1+\bar x)}{\bar x}=\frac{1}{\bar x}-\bar x$、$\frac{\bar y(1+y)}{y}=\frac1y-y$ 后各给出 $6\pi\as r_\chi(2X_A^2-X_A)$。可见“常数项”只在为发散积分规定数值之后才有意义（脚本 \texttt{s2\_\allowbreak qcdf.py}、\texttt{s8\_\allowbreak review.py}）。
\end{proof}

\paragraph{地位与缺口}
破坏（[V]），证据为显式计算：在标准共线因子化中这些幂次修正的卷积发散。$X_{A,H}$ 只是对未知非微扰贡献的参数化：BBNS 称对所有湮灭项取同一个 $X_A$“is a crude model”（是一个粗糙的模型）\cite[第~3.5~节]{BenekeEtAl2001}，Beneke 与 Neubert 称弱湮灭的处理“is model-dependent”（依赖于模型）\cite[第~2.4~节]{BenekeNeubert2003}。两文把“默认值”印作 $\ln(m_M/\Lambda_h)\approx2.4$，与所用参数化不符，应为 $\ln(m_B/\Lambda_h)=2.36$（附录~Q 第~7~条）。把这些端点发散再因子化成有限的因子化定理，对本过程是猜想性质的（[C]）。其他过程中已有逐例的成功：$h\to\gamma\gamma$（经 $b$ 夸克圈）\cite{Liu:2019oav,Liu:2020wbn}、$gg\to h$（经轻夸克圈）\cite{Liu:2022ajh}、$\bar B\to X_s\gamma$ 的次领头 $O_8$--$O_8$ 贡献\cite{Hurth:2023paz}。对纯弱湮灭衰变，据该文摘要，Lü 等人指出一个此前未计入的领头幂次硬共线贡献，认为它有助于理解卷积中的快度发散\cite{Lu:2022kos}；对纯湮灭道 $\bar B^0\to K^+K^-$，一篇 2024 年的博士论文（未在期刊发表）在树图匹配层次上论证了端点发散在标准贡献与新出现的贡献之间相消\cite{Stillger:2024sgy}。色压低的 $\bar B^0\to D^0M^0$ 中弱湮灭没有幂次压低\cite{Mantry:2003uz}。据我们截至 2026 年 10 月的 INSPIRE 检索，对一般的弱湮灭项与扭度 3 硬旁观者项尚无经同行评议的全阶再因子化；一般情形见第~21~章的\cref{prop:power-nlp-endpoint}（[P]）与\cref{sec:front-nlp}。解决途径是为本过程建立再因子化，那将是一个形式不同的新定理。
